\documentclass[11pt,a4paper]{article}
\usepackage[margin=2.3cm]{geometry}
\usepackage{amsmath,amssymb,mathtools}
\usepackage{xcolor}
\usepackage{booktabs}
\usepackage{enumitem}
\usepackage[most]{tcolorbox}
\usepackage{hyperref}

\newcommand{\R}{\mathbb{R}}
\newcommand{\ind}{\mathbf{1}}
\newcommand{\relu}{\operatorname{ReLU}}
\newcommand{\Diag}{\operatorname{Diag}}
% Colour convention
\definecolor{chg}{RGB}{200,0,0}
\newcommand{\C}[1]{{\color{chg}#1}}          % change / addition (red)
\newcommand{\fix}[2]{#1\;\C{\to\;#2}}         % your expression -> corrected one

\tcbset{boxstyle/.style={colback=gray!5,colframe=gray!60,boxrule=0.5pt,arc=2pt,left=6pt,right=6pt,top=4pt,bottom=4pt}}

\title{MMC201 -- Review of Question 7\\[2pt]\large One gradient-descent step for a one-hidden-layer network}
\author{Gustavo Lozano}
\date{\today}

\begin{document}
\maketitle

\begin{tcolorbox}[boxstyle]
\textbf{Sources.} Problem statement: \texttt{11-nn-exercises.pdf}, Question~7.
Hand derivation: \texttt{11-nn.rnote}, pp.~10--14 (source of the black content in this document).

\medskip
\textbf{Colour legend.} Black: your derivation, transcribed from \texttt{11-nn.rnote} (pp.~10--14), and
the problem statement. \C{Red: my changes} --- shape corrections, the compact $\delta$ notation, the numerical
values and the general rule. Notation $a\;\C{\to\;b}$ means ``you wrote $a$, I replace it by $b$''.
\end{tcolorbox}

\section{Problem}
Statement taken from \texttt{11-nn-exercises.pdf}, Question~7.

\[
Z_1 = W_1X + B_1,\qquad H = \relu(Z_1),\qquad Y = W_2H + B_2,\qquad
E = \tfrac12\lVert Y - Y_*\rVert^2,
\]
\[
X=\begin{pmatrix}1\\2\end{pmatrix},\;
W_1=\begin{pmatrix}1&-1\\0&1\end{pmatrix},\;
B_1=\begin{pmatrix}1\\-1\end{pmatrix},\;
W_2=\begin{pmatrix}2&-1\end{pmatrix},\;
B_2=(0),\;
Y_*=(1),\;
\eta=0.1 .
\]
Shapes: $X,B_1,Z_1,H\in\R^{2\times1}$, $W_1\in\R^{2\times2}$, $W_2\in\R^{1\times2}$, $Y\in\R^{1\times1}$.

\C{\textbf{Convention.} Every gradient $\nabla_\theta E$ has the same shape as $\theta$, so that
$\theta_{\text{nuevo}}=\theta-\eta\nabla_\theta E$ is shape-compatible. $\relu'(0)=0$ (as in Question~1).}

\section{Setup}

\[
\theta=\mathrm{FLAT}(W_1,B_1,W_2,B_2),\qquad
\theta_{k+1}=\theta_k+\alpha_k p_k,\quad p_k=-\nabla_\theta E(\theta_k),\quad \alpha_k=\eta=0.1 .
\]
\C{Per parameter: \quad $W_{\ell,\text{nuevo}}=W_\ell-\eta\nabla_{W_\ell}E,\qquad B_{\ell,\text{nuevo}}=B_\ell-\eta\nabla_{B_\ell}E,\qquad \ell=1,2.$}

\section{Derivation}

\subsection{$\nabla_{B_2}E$}
Chain $B_2\to Y\to E$:
\[
\frac{\partial E}{\partial B_2}=\frac{\partial E}{\partial Y}\cdot\frac{\partial Y}{\partial B_2},\qquad
\frac{\partial E}{\partial Y}=Y-Y_*\in\R^{1\times1}\;(1),\qquad
\frac{\partial Y}{\partial B_2}=1 .
\]
\[
\frac{\partial E}{\partial B_2}=Y-Y_*\quad(2)\qquad \C{\eqqcolon\;\delta_2 .}
\]
\C{$\delta_2=\nabla_YE$ is the \emph{output error signal}; it is reused in every other gradient.}

\subsection{$\nabla_{W_2}E$}
Chain $W_2\to Y\to E$:
\[
\nabla_{W_2}E=\Bigl(\frac{\partial E}{\partial w_{2,1j}}\Bigr)_{j=1,2}\in\fix{\R^{2\times1}}{\R^{1\times2}},
\qquad
\frac{\partial Y}{\partial w_{2,1j}}=\frac{\partial}{\partial w_{2,1j}}\Bigl(\sum_{m=1}^{2}w_{2,1m}H_m+B_2\Bigr)=H_j,
\]
\[
\frac{\partial E}{\partial w_{2,1j}}=\frac{\partial E}{\partial Y}\,H_j,
\qquad
\nabla_{W_2}E=\fix{\frac{\partial E}{\partial Y}\,H\in\R^{2\times1}\;(3)}{\;\delta_2\,H^{\top}\in\R^{1\times2}} .
\]
\C{Reason: $W_2$ is $1\times2$, so its gradient must be a row. Entries unchanged.}

\subsection{$\nabla_{B_1}E$}
Chain $B_1\to Z_1\to H\to Y\to E$:
\[
\frac{\partial E}{\partial B_{1,j}}
=\underbrace{\frac{\partial E}{\partial Y}}_{1\times1}\,
 \underbrace{\frac{\partial Y}{\partial H}}_{1\times2}\,
 \underbrace{\frac{\partial H}{\partial Z_1}}_{2\times2}\,
 \underbrace{\frac{\partial Z_1}{\partial B_{1,j}}}_{2\times1},
\]
\[
\frac{\partial Y}{\partial H}=W_2\in\R^{1\times2}\;(4),\qquad
\frac{\partial H_i}{\partial z_{1,j}}=\delta_{ij}\,\ind\{z_{1,j}>0\}\;(5),\qquad
\frac{\partial z_{1,i}}{\partial B_{1,j}}=\delta_{ij}.
\]
\[
\fix{\nabla_{B_1}E}{\dfrac{\partial E}{\partial B_1}}=\begin{bmatrix}\dfrac{\partial E}{\partial B_{1,1}}\\[8pt]\dfrac{\partial E}{\partial B_{1,2}}\end{bmatrix}^{\!\top}
=(Y-Y_*)\,W_2\begin{bmatrix}\ind\{z_{1,1}>0\}&0\\0&\ind\{z_{1,2}>0\}\end{bmatrix}\in\R^{1\times2}
\]
{\color{chg}
Let $r\coloneqq\relu'(Z_1)=\begin{pmatrix}\ind\{z_{1,1}>0\}\\\ind\{z_{1,2}>0\}\end{pmatrix}$, so that
$\begin{bmatrix}\ind\{z_{1,1}>0\}&0\\0&\ind\{z_{1,2}>0\}\end{bmatrix}=\Diag(r)$, and $\delta_2\coloneqq Y-Y_*$.
\begin{align*}
\nabla_{B_1}E
&\coloneqq\Bigl(\frac{\partial E}{\partial B_1}\Bigr)^{\top}\in\R^{2\times1}
&&\\
&=\Bigl[(Y-Y_*)\,W_2\,\Diag(r)\Bigr]^{\top}
&&\\
&=\Diag(r)^{\top}\,W_2^{\top}\,(Y-Y_*)^{\top}
&&(ABC)^{\top}=C^{\top}B^{\top}A^{\top}\\
&=\Diag(r)\,W_2^{\top}\,\delta_2
&&\Diag(r)^{\top}=\Diag(r),\;\;\delta_2\in\R^{1\times1}\Rightarrow\delta_2^{\top}=\delta_2\\
&=\begin{bmatrix}r_1&0\\0&r_2\end{bmatrix}
  \begin{bmatrix}(W_2^{\top}\delta_2)_1\\(W_2^{\top}\delta_2)_2\end{bmatrix}
&&\\
&=\begin{bmatrix}(W_2^{\top}\delta_2)_1\,r_1\\(W_2^{\top}\delta_2)_2\,r_2\end{bmatrix}
&&\\
&=\bigl(W_2^{\top}\delta_2\bigr)\odot r
&&(u\odot v)_i=u_i\,v_i\\
&=\bigl(W_2^{\top}\delta_2\bigr)\odot\relu'(Z_1)\;\eqqcolon\;\delta_1\in\R^{2\times1}.
&&
\end{align*}
\[
\frac{\partial Z_1}{\partial B_1}=I_2\;\Longrightarrow\;\nabla_{B_1}E=\nabla_{Z_1}E=\delta_1 .
\]
}

\subsection{$\nabla_{W_1}E$}
Chain $W_1\to Z_1\to H\to Y\to E$:
\[
\nabla_{W_1}E=\Bigl(\frac{\partial E}{\partial w_{1,ij}}\Bigr)_{i,j=1,2}\in\fix{\R^{4\times1}}{\R^{2\times2}},
\qquad
\frac{\partial E}{\partial w_{1,ij}}
=\underbrace{\frac{\partial E}{\partial Y}}_{\fix{1\times2}{1\times1}}\,
 \underbrace{\frac{\partial Y}{\partial H}}_{1\times2}\,
 \underbrace{\frac{\partial H}{\partial Z_1}}_{2\times2}\,
 \underbrace{\frac{\partial Z_1}{\partial w_{1,ij}}}_{2\times1},
\]
\[
\frac{\partial z_{1,k}}{\partial w_{1,ij}}=\frac{\partial}{\partial w_{1,ij}}\Bigl(\sum_{m=1}^{2}w_{1,km}x_m+B_{1,k}\Bigr)=\delta_{ki}\,x_j,
\qquad
\frac{\partial Z_1}{\partial w_{1,ij}}=\begin{bmatrix}\delta_{1i}\,x_j\\\delta_{2i}\,x_j\end{bmatrix}\in\R^{2\times1},
\]
\[
\frac{\partial E}{\partial w_{1,ij}}=(Y-Y_*)\,W_2\begin{bmatrix}\ind\{z_{1,1}>0\}&0\\0&\ind\{z_{1,2}>0\}\end{bmatrix}
\begin{bmatrix}\delta_{1i}\,x_j\\\delta_{2i}\,x_j\end{bmatrix}
\]
{\color{chg}
With $r=\relu'(Z_1)$, $\delta_2=Y-Y_*$ and $\delta_1=\Diag(r)\,W_2^{\top}\delta_2$ (Section~3.3):
\begin{align*}
\frac{\partial E}{\partial w_{1,ij}}
&=(Y-Y_*)\,W_2\,\Diag(r)\begin{bmatrix}\delta_{1i}\,x_j\\\delta_{2i}\,x_j\end{bmatrix}
&&\\
&=\bigl(\Diag(r)\,W_2^{\top}\,\delta_2\bigr)^{\top}\begin{bmatrix}\delta_{1i}\,x_j\\\delta_{2i}\,x_j\end{bmatrix}
&&A^{\top}B^{\top}C^{\top}=(CBA)^{\top},\;\Diag(r)^{\top}=\Diag(r),\;\delta_2^{\top}=\delta_2\\
&=\delta_1^{\top}\begin{bmatrix}\delta_{1i}\,x_j\\\delta_{2i}\,x_j\end{bmatrix}
&&\\
&=\begin{bmatrix}(\delta_1)_1&(\delta_1)_2\end{bmatrix}\begin{bmatrix}\delta_{1i}\,x_j\\\delta_{2i}\,x_j\end{bmatrix}
&&\\
&=(\delta_1)_1\,\delta_{1i}\,x_j+(\delta_1)_2\,\delta_{2i}\,x_j
&&\\
&=\sum_{k=1}^{2}(\delta_1)_k\,\delta_{ki}\,x_j
&&\\
&=(\delta_1)_i\,x_j
&&\textstyle\sum_k a_k\,\delta_{ki}=a_i\\
&=\bigl(\delta_1X^{\top}\bigr)_{ij}
&&(uv^{\top})_{ij}=u_i\,v_j .
\end{align*}
}
\[
\nabla_{W_1}E=\begin{bmatrix}
(Y-Y_*)\,W_2\begin{bmatrix}\ind\{z_{1,1}>0\}\,x_1\\0\end{bmatrix}\\[10pt]
(Y-Y_*)\,W_2\begin{bmatrix}\ind\{z_{1,1}>0\}\,x_2\\0\end{bmatrix}\\[10pt]
(Y-Y_*)\,W_2\begin{bmatrix}0\\\ind\{z_{1,2}>0\}\,x_1\end{bmatrix}\\[10pt]
(Y-Y_*)\,W_2\begin{bmatrix}0\\\ind\{z_{1,2}>0\}\,x_2\end{bmatrix}
\end{bmatrix}
\;\C{\to\;\begin{bmatrix}\dfrac{\partial E}{\partial w_{1,11}}&\dfrac{\partial E}{\partial w_{1,12}}\\[8pt]
                 \dfrac{\partial E}{\partial w_{1,21}}&\dfrac{\partial E}{\partial w_{1,22}}\end{bmatrix}\in\R^{2\times2}}
\]
{\color{chg}
\begin{align*}
\nabla_{W_1}E
&=\begin{bmatrix}(\delta_1)_1x_1&(\delta_1)_1x_2\\(\delta_1)_2x_1&(\delta_1)_2x_2\end{bmatrix}
&&\frac{\partial E}{\partial w_{1,ij}}=(\delta_1)_i\,x_j\\
&=\begin{bmatrix}(\delta_1)_1\\(\delta_1)_2\end{bmatrix}\begin{bmatrix}x_1&x_2\end{bmatrix}
&&(uv^{\top})_{ij}=u_i\,v_j\\
&=\delta_1X^{\top}\in\R^{2\times2}.
&&
\end{align*}
}

\section{\C{Summary in compact form}}

{\color{chg}
\begin{tcolorbox}[boxstyle]
\[
\delta_2=Y-Y_*,\qquad \delta_1=\bigl(W_2^{\top}\delta_2\bigr)\odot\relu'(Z_1),
\]
\[
\nabla_{W_2}E=\delta_2H^{\top},\quad \nabla_{B_2}E=\delta_2,\qquad
\nabla_{W_1}E=\delta_1X^{\top},\quad \nabla_{B_1}E=\delta_1 .
\]
\end{tcolorbox}

\begin{center}
\renewcommand{\arraystretch}{1.35}
\begin{tabular}{@{}lll@{}}
\toprule
Gradient & {\color{black}Your notes} & Compact form \\
\midrule
$\nabla_{B_2}E$ & {\color{black}$Y-Y_*$} & $\delta_2$ \\
$\nabla_{W_2}E$ & {\color{black}$\frac{\partial E}{\partial Y}H\in\R^{2\times1}$} & $\delta_2H^{\top}\in\R^{1\times2}$ \\
$\nabla_{B_1}E$ & {\color{black}$(Y-Y_*)\,W_2\,\Diag(\ind\{Z_1>0\})\in\R^{1\times2}$} & $\delta_1\in\R^{2\times1}$ \\
$\nabla_{W_1}E$ & {\color{black}four blocks, $\R^{4\times1}$} & $\delta_1X^{\top}\in\R^{2\times2}$ \\
\bottomrule
\end{tabular}
\end{center}
}

\section{\C{Numerical computation}}
{\color{chg}
\paragraph{Forward pass.}
\begin{align*}
Z_1&=W_1X+B_1=\begin{pmatrix}1&-1\\0&1\end{pmatrix}\begin{pmatrix}1\\2\end{pmatrix}+\begin{pmatrix}1\\-1\end{pmatrix}
=\begin{pmatrix}-1\\2\end{pmatrix}+\begin{pmatrix}1\\-1\end{pmatrix}=\begin{pmatrix}0\\1\end{pmatrix},\\
H&=\relu(Z_1)=\begin{pmatrix}\max(0,0)\\\max(0,1)\end{pmatrix}=\begin{pmatrix}0\\1\end{pmatrix},\\
Y&=W_2H+B_2=\begin{pmatrix}2&-1\end{pmatrix}\begin{pmatrix}0\\1\end{pmatrix}+0=-1,\\
E&=\tfrac12\lVert Y-Y_*\rVert^2=\tfrac12(-1-1)^2=2.
\end{align*}

\paragraph{Backward pass.}
\[
\delta_2=-2,\qquad
W_2^{\top}\delta_2=\begin{pmatrix}-4\\2\end{pmatrix},\qquad
\relu'(Z_1)=\begin{pmatrix}0\\1\end{pmatrix},\qquad
\delta_1=\begin{pmatrix}-4\\2\end{pmatrix}\odot\begin{pmatrix}0\\1\end{pmatrix}=\begin{pmatrix}0\\2\end{pmatrix}.
\]
\emph{Caution:} $z_{1,1}=0$ lies exactly on the ReLU kink. With $\relu'(0)=0$ the first hidden unit passes
no gradient. Using $\relu'(0)=1$ would give $\delta_1=(-4,2)^{\top}$ and different updates.

\paragraph{Gradients.}
\[
\nabla_{B_2}E=-2,\quad
\nabla_{W_2}E=\begin{pmatrix}0&-2\end{pmatrix},\quad
\nabla_{B_1}E=\begin{pmatrix}0\\2\end{pmatrix},\quad
\nabla_{W_1}E=\begin{pmatrix}0\\2\end{pmatrix}\begin{pmatrix}1&2\end{pmatrix}=\begin{pmatrix}0&0\\2&4\end{pmatrix}.
\]

\paragraph{Updates ($\eta=0.1$).}
\[
W_{1,\text{nuevo}}=\begin{pmatrix}1&-1\\-0.2&0.6\end{pmatrix},\quad
B_{1,\text{nuevo}}=\begin{pmatrix}1\\-1.2\end{pmatrix},\quad
W_{2,\text{nuevo}}=\begin{pmatrix}2&-0.8\end{pmatrix},\quad
B_{2,\text{nuevo}}=0.2 .
\]

\paragraph{Sanity check.}
With the new parameters, $Z_1=(0,-0.2)^{\top}$, $H=(0,0)^{\top}$, $Y=0.2$, $E=\tfrac12(0.8)^2=0.32<2$.
}

\section{\C{General rule}}
{\color{chg}
For $Z_\ell=W_\ell A_{\ell-1}+B_\ell$, $A_\ell=f_\ell(Z_\ell)$, $A_0=X$, $\ell=1,\dots,L$:
\[
\delta_L=\nabla_{A_L}E\odot f_L'(Z_L),\qquad
\delta_\ell=\bigl(W_{\ell+1}^{\top}\delta_{\ell+1}\bigr)\odot f_\ell'(Z_\ell),\qquad
\nabla_{W_\ell}E=\delta_\ell A_{\ell-1}^{\top},\qquad \nabla_{B_\ell}E=\delta_\ell .
\]
Question 7 is $L=2$, $f_1=\relu$, $f_2=\mathrm{id}$, $A_1=H$. Questions 4--6 are its single-layer pieces.
}

\end{document}
